The interventions in this appendix are diagnostic controls rather than proposed
universal fixes. Each intervention removes or stabilizes one geometry channel so
that detector behavior can be interpreted more mechanistically.

\subsection{L2 Normalization}
\label{app:l2-normalization}

Feature normalization removes radial scale:
\[
    \hat{f}_\theta(x)
    =
    \frac{f_\theta(x)}
    {\|f_\theta(x)\|_2+\epsilon}.
\]
Raw and L2-normalized detectors are compared under the same optimizer, dataset,
OOD split, and seed. If raw Mahalanobis or raw kNN performs poorly for one
optimizer but recovers after L2 normalization, the failure includes a
norm-driven component. If a gap remains after L2 normalization, radial scale is
not sufficient to explain the detector behavior.

\subsection{Covariance Shrinkage}
\label{app:covariance-shrinkage}

For class covariance \(\Sigma_c\), shrinkage uses
\[
    \Sigma_{c,\alpha}
    =
    (1-\alpha)\Sigma_c
    +
    \alpha\frac{\operatorname{tr}(\Sigma_c)}{d}I.
\]
This stabilizes inverse covariance by increasing small eigenvalues and pulling
large eigenvalues toward the mean. Tied, diagonal, and shrinkage covariance
variants are treated as diagnostic controls. If GMM or Mahalanobis scores
recover after shrinkage, covariance estimation or covariance conditioning
contributes to the detector behavior.

\subsection{PCA Projection and Residual Diagnostics}
\label{app:pca-projection-residual}

PCA projection preserves the ID principal subspace:
\[
    \hat{f}_\theta(x)
    =
    PP^\top(f_\theta(x)-\mu)+\mu.
\]
A PCA reconstruction-error style score reads the residual component:
\[
    S_{\mathrm{PCARec}}(x)
    =
    -\|(I-PP^\top)(f_\theta(x)-\mu)\|_2.
\]
The negative sign follows the higher-is-ID-like convention. These diagnostics
test whether optimizer-induced geometry differences concentrate inside the ID
principal subspace or in residual directions. The PCA fitting split, retained
dimension rule, and score direction should be fixed before reporting PCA or
NeCo scores as main results.

\subsection{Angular Alignment Probes}
\label{app:angular-alignment-probes}

CTM is an angular alignment probe because it removes feature norm and reads
cosine similarity between classifier weights and features. It should be compared
with NC3 classifier-mean alignment rather than interpreted as a generic
feature-distance score. CTM is reported in the main text only when the official
score convention and implementation are fixed for the experiment package.

\subsection{Implementation Notes and Score Convention}
\label{app:diagnostic-implementation-notes}

Diagnostic controls should use the same training checkpoints, feature extraction
layer, ID split, OOD datasets, and score direction convention as the corresponding
raw detector. Manual post-hoc variants should be labeled explicitly as
diagnostics. They should not be presented as faithful reproductions of methods
whose original training recipe includes additional regularizers or architectural
constraints.
